% ============================================================================
% 🛑 INCOMPLETO -- NAO ENTREGAR.  Nao esta em `plates/`, e nao deve entrar sem
%    uma rodada de reposicionamento manual.  A chapa VIVA e' `hgi_flow.tex`.
%
% O QUE ISTO E'.  Tentativa (27/08) de saldar a divida declarada pela `ppt`: o
% slide do HGI usa `\includegraphics[width=0.70\textwidth]`, o que AFINA OS
% TRACOS em 30%.  O certo seria regenerar a chapa no tamanho final, com as
% espessuras recalculadas em vez de escaladas.
%
% O QUE FOI FEITO.  Transformacao mecanica: `scale=0.67` nas coordenadas, todas
% as dimensoes absolutas (minimum *, text width, inner sep, xshift/yshift,
% shorten, rounded corners, raio dos circulos do pictograma) x0,70, fonte fixada
% em 6,7pt -- e `line width` DELIBERADAMENTE INTOCADO, que era o objetivo todo.
% Saida: 257,57 x 113,82 pt (contra 264,5 x 115,0 renderizados hoje) -- MENOR
% nas duas dimensoes, portanto seguro para o frame.
%
% ⚠ POR QUE NAO SERVE, e o resultado negativo vale mais que a tentativa:
% UMA TRANSFORMACAO MECANICA NAO REGENERA ESTA CHAPA.  Duas coisas quebram:
%   1. os PICTOGRAMAS nao encolhem na mesma proporcao das caixas -- as caixas
%      seguem `minimum height`/`text width` (escalados x0,70) e o texto a 6,7pt,
%      mas os pentagonos seguem so' as coordenadas; o resultado e' que eles
%      ficam ~17% maiores em relacao ao resto;
%   2. os rotulos `Delaunay` e `Geographic adjacency` foram posicionados A MAO
%      contra a geometria antiga, e ENCOSTAM nos grafos depois da transformacao.
%      `Geographic adjacency` e' o pior: invade o pictograma e a caixa
%      `Region GCN`.
%
% O QUE FALTA, para quem retomar: reposicionar esses dois rotulos e reequilibrar
% o tamanho do `gpic` contra as caixas -- ou seja, a mesma rodada de composicao
% que a chapa original custou.  Nao e' trabalho de `sed`.
%
% ⚠ NAO ADIANTA MEXER SO' NO `scale` PARA FECHAR A PROPORCAO.  Varri
% (scale, fonte) = (0,70;7,0) (0,70;6,6) (0,68;6,8) (0,68;6,6) (0,66;6,6): a
% razao largura/altura fica presa em ~2,26 e NUNCA chega aos 2,299 do original.
% E' estrutural -- os tracos em peso cheio SAO a altura extra.  Nao existe
% versao com traco cheio E a mesma altura renderizada; so' da' para escolher de
% que lado fica o erro.  (Por isso 0,67: fica menor nas duas dimensoes.)
% ============================================================================

% ============================================================================
% HGI (Hierarchical Graph Infomax) -- diagrama de slide (impresso 23).
% Fonte: `presentation/hgi_draw.txt` (o ASCII do autor) + a revisao de
% composicao que ele fez na rodada 7.
%
% POR QUE DUAS FAIXAS, e nao uma linha so'.  O fluxo forward pedido tem ONZE
% estagios (cat.encoder, grafos POI, POI GCN, vetores, MHA, vetores, grafos de
% regiao, Region GCN, vetores, agregacao por area, cidade).  Somadas, as caixas
% dao 11,8cm; sobram ~2mm por vao numa largura util de 14cm, o que produz setas
% menores que a propria ponta.  O autor pediu "uma nova linha" para os grafos de
% regiao, e e' exatamente isso que resolve: faixa 1 ate' os region raw
% embeddings, faixa 2 dali ate' a cidade.
% A faixa 2 e' alinhada A' DIREITA de proposito: assim a volta de uma faixa para
% a outra e' curta (1,2cm) em vez de atravessar a figura inteira -- e o canto
% inferior esquerdo que sobra e' onde a banda de discriminacao cabe.
%
% GRAMATICA HERDADA DA CHAPA DO DGI (as tres chapas sao familia):
%   VERMELHO = ramo negativo, e nada mais.  Grafo corrompido, vetores dele e as
%   setas que os carregam.  Por isso o backprop e' CINZA tracejado: se vermelho
%   significa "negativo", nao pode significar tambem "sinal de treino".
%
% POR QUE AS LIGACOES DE DISCRIMINACAO SAO TRACEJADAS: elas atravessam o fluxo
% forward (a banda de discriminacao esta' embaixo e as fontes estao nas duas
% faixas).  O tracejado separa "o que flui" de "o que e' comparado".
%
% PRECISOES do hgi_draw.txt que o desenho respeita:
%  - DUAS fronteiras, nunca tres.  Tres seria desenhar o Check2HGI com outro
%    nome e destruiria o contraste que o slide 31 precisa fazer;
%  - so' o ramo ORIGINAL sobe para a cidade; o corrompido para nas region
%    embeddings, onde o discriminador o consome.  No desenho isso e' visivel
%    pela ausencia de seta: a raia vermelha simplesmente nao continua;
%  - a GCN de nivel de regiao existe e a prosa do Cap. 2 a omite.
%
% ⚠ RESSALVA REGISTRADA (rodada 7).  O autor pediu que o "Discrimination
%   POI-Region" receba os DOIS vetores que saem do POI GCN -- o original e o
%   corrompido.  Desenhado com o ORIGINAL apenas, e a razao esta' no proprio
%   hgi_draw.txt dele: "O negativo da fronteira POI-regiao e' um POI DIFERENTE
%   pareado com a MESMA regiao: (p_j, r_k)."  O ramo corrompido serve a'
%   fronteira REGIAO-CIDADE, via r~_k -- e la' ele esta' desenhado.  Ligar o
%   vermelho ao D_pr diria que o negativo daquela fronteira e' um POI com
%   features embaralhadas, que e' outro metodo.  Se ele reafirmar, trocar.
%
% Gabarito: 140 x 63,7 mm.  ⚠ esta chapa usa ~60mm de altura -- e' um slide de
% FIGURA, sem espaco para as duas linhas de texto que o `gate` orcou.
% ============================================================================
% UMA COR POR DISCRIMINACAO.  Nenhuma das duas esta' na paleta do nesped de
% proposito: teal, navy e magenta ja' significam contexto, fluxo e negativo nas
% tres chapas, e reusar qualquer um deles era o que fazia as ligacoes de
% discriminacao se confundirem com o resto do desenho.
%  ambar  = fronteira POI-regiao
%  violeta = fronteira regiao-cidade
% Cada cor vai da FONTE ate' a Total loss sem interrupcao, entao da' para seguir
% uma comparacao inteira com o olho sem soltar o fio.
\definecolor{discA}{rgb}{0.74,0.50,0.05}
\definecolor{discB}{rgb}{0.36,0.29,0.58}
\tikzset{
  gpic/.pic={
    \draw[line width=0.8pt, opacity=0.45]
      (-0.34,0.15)--(-0.06,0.28)--(0.28,0.06)--(0.06,-0.24)--(-0.28,-0.17)--cycle
      (-0.34,0.15)--(0.28,0.06);
    \fill (-0.34,0.15) circle (0.91pt) (-0.06,0.28) circle (0.91pt)
          (0.28,0.06) circle (0.91pt) (0.06,-0.24) circle (0.91pt)
          (-0.28,-0.17) circle (0.91pt);
  }
}
\begin{tikzpicture}[
    scale=0.67,
    font=\fontsize{6.7}{8.0}\selectfont,
    >={Latex[length=1.4mm]},
    lbl/.style={align=center, font=\fontsize{6.7}{8.0}\selectfont, text=secondary},
    box/.style={draw=secondaryshade, rounded corners=1.75pt, fill=white,
                align=center, inner sep=2.8pt, line width=1pt, text=secondary},
    tall/.style={box, minimum height=0.945cm},
    dboxA/.style={box, draw=discA, line width=1.1pt, minimum height=0.525cm},
    dboxB/.style={box, draw=discB, line width=1.1pt, minimum height=0.525cm},
    dboxL/.style={box, draw=black!55, line width=1.1pt, minimum height=0.525cm},
    vcell/.style={draw=secondary!70, line width=0.8pt, minimum width=2.38mm,
                  minimum height=1.26mm, inner sep=0pt, fill=white},
    ncell/.style={vcell, draw=alert!85, fill=alert!10},
    ccell/.style={vcell, draw=primaryshade, fill=primarytint},
    conn/.style={circle, draw=secondary, line width=1pt, fill=white,
                 minimum size=3.5mm, inner sep=0pt, font=\fontsize{6.7}{8.0}\selectfont, text=secondary},
    nconn/.style={conn, draw=alert!85, text=alert},
    flow/.style={->, line width=1.1pt, draw=secondary},
    nflow/.style={->, line width=1.1pt, draw=alert!85},
    cmpA/.style={->, line width=0.95pt, draw=discA, dashed, shorten >=1.4pt},
    cmpB/.style={->, line width=0.95pt, draw=discB, dashed, shorten >=1.4pt}
]
\def\yA{2.30}\def\yB{1.25}\def\yC{0.20}\def\yD{-0.85}

% ======================= LINHA 1 -- nivel POI, do encoder ao region raw
\node[tall, text width=1.12cm] (cat) at (1.10,1.775)
      {\textbf{Category}\\\textbf{encoder}};
{\color{secondary}\pic at (3.30,\yA) {gpic};}
{\color{alert}\pic at (3.30,\yB) {gpic};}
\node[lbl] at (3.30,1.775) {\textbf{Delaunay}};
\draw[flow]  ([yshift=3.675mm]cat.east)  -- (2.85,\yA);
\draw[nflow] ([yshift=-3.675mm]cat.east) -- (2.85,\yB);

\node[tall, text width=0.77cm] (pgcn) at (5.10,1.775) {\textbf{POI}\\\textbf{GCN}};
\draw[flow]  (3.75,\yA) -- ([yshift=3.675mm]pgcn.west);
\draw[nflow] (3.75,\yB) -- ([yshift=-3.675mm]pgcn.west);

\begin{scope}[local bounding box=pva]
  \foreach \k in {-1,0,1} \node[vcell] at (6.80,\yA+\k*0.18) {};\end{scope}
\begin{scope}[local bounding box=pvb]
  \foreach \k in {-1,0,1} \node[ncell] at (6.80,\yB+\k*0.18) {};\end{scope}
\draw[flow]  ([yshift=3.675mm]pgcn.east)  -- (pva.west);
\draw[nflow] ([yshift=-3.675mm]pgcn.east) -- (pvb.west);

\node[tall, text width=1.085cm] (mha) at (8.95,1.775)
      {\textbf{Multi-Head}\\\textbf{Attention}};
\draw[flow]  (pva.east) -- ([yshift=3.675mm]mha.west);
\draw[nflow] (pvb.east) -- ([yshift=-3.675mm]mha.west);

\begin{scope}[local bounding box=rra]
  \foreach \k in {-1,0,1} \node[vcell] at (10.85,\yA+\k*0.18) {};\end{scope}
\begin{scope}[local bounding box=rrb]
  \foreach \k in {-1,0,1} \node[ncell] at (10.85,\yB+\k*0.18) {};\end{scope}
\draw[flow]  ([yshift=3.675mm]mha.east)  -- (rra.west);
\draw[nflow] ([yshift=-3.675mm]mha.east) -- (rrb.west);

\node[conn]  (ca1) at (11.95,\yA) {A};
\node[nconn] (cb1) at (11.95,\yB) {B};
\draw[flow]  (rra.east) -- (ca1.west);
\draw[nflow] (rrb.east) -- (cb1.west);

% ======================= LINHA 2 -- nivel regiao, do grafo a' cidade
\node[conn]  (ca2) at (0.75,\yC) {A};
\node[nconn] (cb2) at (0.75,\yD) {B};
{\color{secondary}\pic at (1.95,\yC) {gpic};}
{\color{alert}\pic at (1.95,\yD) {gpic};}
\node[lbl] at (1.95,-0.325) {\textbf{Geographic adjacency}};
\draw[flow]  (ca2.east) -- (1.50,\yC);
\draw[nflow] (cb2.east) -- (1.50,\yD);

\node[tall, text width=0.875cm] (rgcn) at (4.60,-0.325) {\textbf{Region}\\\textbf{GCN}};
\draw[flow]  (2.40,\yC) -- ([yshift=3.675mm]rgcn.west);
\draw[nflow] (2.40,\yD) -- ([yshift=-3.675mm]rgcn.west);

\begin{scope}[local bounding box=rea]
  \foreach \k in {-1,0,1} \node[vcell] at (5.85,\yC+\k*0.18) {};\end{scope}
\begin{scope}[local bounding box=reb]
  \foreach \k in {-1,0,1} \node[ncell] at (5.85,\yD+\k*0.18) {};\end{scope}
\draw[flow]  ([yshift=3.675mm]rgcn.east)  -- (rea.west);
\draw[nflow] ([yshift=-3.675mm]rgcn.east) -- (reb.west);

% so' o ramo ORIGINAL sobe: a raia vermelha simplesmente nao continua daqui.
\node[tall, text width=1.295cm, minimum height=0.7cm] (awa) at (8.10,\yC)
      {\textbf{Area-Weighted}\\\textbf{Aggregation}};
\draw[flow] (rea.east) -- (awa.west);
\begin{scope}[local bounding box=city]
  \foreach \k in {-1,0,1} \node[ccell] at (10.10,\yC+\k*0.18) {};\end{scope}
\draw[flow] (awa.east) -- (city.west);
\node[lbl, text=primaryshade, anchor=west] at (10.40,\yC) {city};

% ======================= LINHA 3 -- so' discriminacao, ate' a Total loss.
% POI-Regiao a' ESQUERDA e Regiao-Cidade a' DIREITA: a ordem da hierarquia
% (POI->regiao vem antes de regiao->cidade) passa a ser tambem a ordem de
% leitura da faixa.
\node[dboxB, text width=1.505cm] (drc) at (5.00,-2.675)
      {Discrimination\\Region--City};
\node[dboxL, text width=0.7cm] (loss) at (9.00,-2.675) {Total\\loss};
\node[dboxA, text width=1.505cm] (dpr) at (12.00,-2.675)
      {Discrimination\\POI--Region};
\draw[cmpB] (drc.east) -- (loss.west);
\draw[cmpA] (dpr.west) -- (loss.east);

% TODA ligacao de comparacao sai pela BASE do embedding, nunca pela lateral.
% Motivo: a lateral e' onde as setas do forward CHEGAM, e uma saida ali fica
% ambigua -- o fio parecia nascer no embedding de baixo, o rosa, quando nasce no
% de cima.  O degrau curto logo apos a saida serve para passar ao largo do
% embedding da outra raia, que fica exatamente abaixo.
%
% AS PONTAS CHEGAM ALINHADAS: duas entradas igualmente espacadas no topo do
% POI-Regiao (4,40 / 5,60) e tres no do Regiao-Cidade (11,30 / 12,00 / 12,70).
%
% QUATRO CANAIS, e a ordem deles nao e' arbitraria.  Com todos os fios indo para
% o mesmo lado, dois da MESMA COR so' deixam de se cruzar se o canal mais alto
% for o que dobra mais longe na direcao do percurso.  Dai a atribuicao abaixo --
% e por isso `pva` e `city` PODEM dividir o canal -1,30: os trechos horizontais
% deles nao se sobrepoem em x (4,40..6,45 contra 10,10..12,70).
% Cruzamentos ambar-com-violeta sobram e sao aceitaveis: a cor resolve.  O que
% nao pode e' violeta cruzar violeta -- ali a cor deixa de ajudar.

% --- fronteira POI-REGIAO (ambar), para a esquerda
\draw[cmpA] (pva.south) -- (6.80,1.80) -- (6.45,1.80) -- (6.45,-1.725)
            -- (12.00,-1.725) -- (12.00,-2.30);
\draw[cmpA] (pvb.south) -- (6.80,-1.875) -- (11.25,-1.875) -- (11.25,-2.30);
\draw[cmpA] (rra.south) -- (10.85,1.80) -- (12.75,1.80) -- (12.75,-2.30);

% --- fronteira REGIAO-CIDADE (violeta), para a direita
\draw[cmpB] (reb.south) -- (5.85,-1.275) -- (4.30,-1.275) -- (4.30,-2.30);
\draw[cmpB] (rea.south) -- (5.85,-0.30) -- (6.20,-0.30) -- (6.20,-1.425)
            -- (5.00,-1.425) -- (5.00,-2.30);
\draw[cmpB] (city.south) -- (10.10,-1.575) -- (5.70,-1.575) -- (5.70,-2.30);

\end{tikzpicture}
